% --- Archon physics formalization source begin ---
% archon:physics
% archon:covers PhyXMiniProblems/problem_phyx_mini_0248.lean
% archon:source-report reports/phyx_mini/problem_phyx_mini_0248.source.json

\chapter{Physics problem phyx\_mini\_0248}
\label{ch:PhyXMiniProblems_problem_phyx_mini_0248}

\paragraph{Problem source.}
\#\# Physical scenario

Light from a helium-neon laser (\$\textbackslash{}lambda = 633\textbackslash{};nm\$) illuminates two slits spaced \$0.40\textbackslash{};mm\$ apart. A viewing screen is \$2.0\textbackslash{};m\$ behind the slits.

\#\# Question

What is the distance between the two \$m = 2\$ dark fringes?

\#\# Answer choices

- A: A: 15.8mm

- B: B: 15.2mm

- C: C: 16.8mm

- D: D: 16.2mm

\#\# Figure caption (auxiliary; use the image as primary evidence)

The image depicts a diagram of a double-slit experiment, illustrating the interference pattern of waves.

1. **Plane Wave**: On the left, parallel purple lines representing a plane wave are incident on a double slit, labeled with a wavelength symbol (λ).

2. **Slits**: Two narrow gaps in a barrier allow the waves to pass through, labeled as the "Top view of the double slit."

3. **Spreading Waves**: Behind the slits, circular purple waves spread out from each slit, creating an interference pattern.

4. **Interference**: The region where the waves overlap shows constructive and destructive interference.

5. **Interference Pattern**: On the right, a screen displays a pattern of bright and dark fringes. The bright fringes occur where the antinodal lines intersect the screen.

6. **Labels and Text**: 
   - On the left: 
     - "1. A plane wave is incident on the double slit."
     - "2. Waves spread out behind each slit."
     - "3. The waves interfere in the region where they overlap."
     - "4. Bright fringes occur where the antinodal lines intersect the viewing screen."
   - On the right, the pattern is labeled:
     - "Central maximum" at m = 0
     - "m = 1, 2, 3, 4" for the subsequent maxima above and below the central maximum.

This diagram explains the basic principles of wave interference in the double-slit experiment.

\#\# Dataset metadata

Wave Properties; Physical Model Grounding Reasoning, Multi-Formula Reasoning

\paragraph{Recorded answer/context.}
A: A: 15.8mm

\paragraph{Figure/image path.}
/root/proposal\_for\_physic/science-mango/phyx\_mini\_run/phyx\_data/test\_image/248.png

\paragraph{Formalization target.}
create a compiling Lean file with sorry bodies at `PhyXMiniProblems/problem\_phyx\_mini\_0248.lean`. The Lean declarations must preserve the physical quantities, dimensions or dimensional roles, figure labels, governing-law hypotheses, and final relation expressed by this problem.
Use Mathlib/Physlib names found through LeanExplore where available. If a physics API is missing, introduce faithful local abstractions rather than scalar placeholder aliases.

\begin{theorem}[Physics formalization target]
\label{thm:physics:phyx_mini_0248:target}
\lean{PhyXMiniProblems.ProblemPhyXMini0248.problem_phyx_mini_0248}
\uses{def:physics:phyx-mini-0248:phyxminiproblems-problemphyxmini0248-doubleslitinterferencesetup, def:physics:phyx-mini-0248:phyxminiproblems-problemphyxmini0248-matchesscenarioandfigure, def:physics:phyx-mini-0248:phyxminiproblems-problemphyxmini0248-matchesproblemreadouts, def:physics:phyx-mini-0248:phyxminiproblems-problemphyxmini0248-hasphysicalparameters, def:physics:phyx-mini-0248:phyxminiproblems-problemphyxmini0248-satisfiesexacttwopointsourcelaws, def:physics:phyx-mini-0248:phyxminiproblems-problemphyxmini0248-darkfringeseparationinmillimeters, def:physics:phyx-mini-0248:phyxminiproblems-problemphyxmini0248-finitescreensymmetricseparation, def:physics:phyx-mini-0248:phyxminiproblems-problemphyxmini0248-recordeddatasetanswer, def:physics:phyx-mini-0248:phyxminiproblems-problemphyxmini0248-isuniquematchingdisplayeddistance}
The assigned autoformalize agent should translate this physics problem into Lean declarations in the covered file, with theorem and lemma proof bodies written as `by sorry`.
\end{theorem}
\begin{proof}
This is an autoformalization task, not a proof task. Produce statements that can later be proved by the physics prover without weakening the physical model.
\end{proof}

% --- Archon declaration topology begin ---
\section{Declaration topology}
\begin{definition}[Length Quantity]
  \label{def:physics:phyx-mini-0248:phyxminiproblems-problemphyxmini0248-lengthquantity}
  \lean{PhyXMiniProblems.ProblemPhyXMini0248.LengthQuantity}
  A nonnegative physical length, independent of a chosen readout unit
\end{definition}
\begin{proof}
  This is the declared model definition of Length Quantity.
\end{proof}
\begin{definition}[Signed Length Quantity]
  \label{def:physics:phyx-mini-0248:phyxminiproblems-problemphyxmini0248-signedlengthquantity}
  \lean{PhyXMiniProblems.ProblemPhyXMini0248.SignedLengthQuantity}
  A signed physical length used for transverse screen coordinates
\end{definition}
\begin{proof}
  This is the declared model definition of Signed Length Quantity.
\end{proof}
\begin{definition}[length Readout]
  \label{def:physics:phyx-mini-0248:phyxminiproblems-problemphyxmini0248-lengthreadout}
  \lean{PhyXMiniProblems.ProblemPhyXMini0248.lengthReadout}
  \uses{def:physics:phyx-mini-0248:phyxminiproblems-problemphyxmini0248-lengthquantity}
  Read a nonnegative physical length in the selected length unit
\end{definition}
\begin{proof}
  This is the declared model definition of length Readout.
\end{proof}
\begin{definition}[signed Length Readout]
  \label{def:physics:phyx-mini-0248:phyxminiproblems-problemphyxmini0248-signedlengthreadout}
  \lean{PhyXMiniProblems.ProblemPhyXMini0248.signedLengthReadout}
  \uses{def:physics:phyx-mini-0248:phyxminiproblems-problemphyxmini0248-signedlengthquantity}
  Read a signed transverse coordinate in the selected length unit
\end{definition}
\begin{proof}
  This is the declared model definition of signed Length Readout.
\end{proof}
\begin{definition}[length In Meters]
  \label{def:physics:phyx-mini-0248:phyxminiproblems-problemphyxmini0248-lengthinmeters}
  \lean{PhyXMiniProblems.ProblemPhyXMini0248.lengthInMeters}
  \uses{def:physics:phyx-mini-0248:phyxminiproblems-problemphyxmini0248-lengthquantity, def:physics:phyx-mini-0248:phyxminiproblems-problemphyxmini0248-lengthreadout}
  Meter readout of a physical length
\end{definition}
\begin{proof}
  This is the declared model definition of length In Meters.
\end{proof}
\begin{definition}[length In Millimeters]
  \label{def:physics:phyx-mini-0248:phyxminiproblems-problemphyxmini0248-lengthinmillimeters}
  \lean{PhyXMiniProblems.ProblemPhyXMini0248.lengthInMillimeters}
  \uses{def:physics:phyx-mini-0248:phyxminiproblems-problemphyxmini0248-lengthquantity, def:physics:phyx-mini-0248:phyxminiproblems-problemphyxmini0248-lengthreadout}
  Millimeter readout of a physical length
\end{definition}
\begin{proof}
  This is the declared model definition of length In Millimeters.
\end{proof}
\begin{definition}[length In Nanometers]
  \label{def:physics:phyx-mini-0248:phyxminiproblems-problemphyxmini0248-lengthinnanometers}
  \lean{PhyXMiniProblems.ProblemPhyXMini0248.lengthInNanometers}
  \uses{def:physics:phyx-mini-0248:phyxminiproblems-problemphyxmini0248-lengthquantity, def:physics:phyx-mini-0248:phyxminiproblems-problemphyxmini0248-lengthreadout}
  Nanometer readout of a physical length
\end{definition}
\begin{proof}
  This is the declared model definition of length In Nanometers.
\end{proof}
\begin{definition}[Slit Label]
  \label{def:physics:phyx-mini-0248:phyxminiproblems-problemphyxmini0248-slitlabel}
  \lean{PhyXMiniProblems.ProblemPhyXMini0248.SlitLabel}
  The two openings in the barrier, viewed from the top as in the figure
\end{definition}
\begin{proof}
  This is the declared model definition of Slit Label.
\end{proof}
\begin{definition}[sign]
  \label{def:physics:phyx-mini-0248:phyxminiproblems-problemphyxmini0248-slitlabel-sign}
  \lean{PhyXMiniProblems.ProblemPhyXMini0248.SlitLabel.sign}
  \uses{def:physics:phyx-mini-0248:phyxminiproblems-problemphyxmini0248-slitlabel}
  The signed transverse location of a slit relative to the optical axis
\end{definition}
\begin{proof}
  This is the declared model definition of sign.
\end{proof}
\begin{definition}[Fringe Side]
  \label{def:physics:phyx-mini-0248:phyxminiproblems-problemphyxmini0248-fringeside}
  \lean{PhyXMiniProblems.ProblemPhyXMini0248.FringeSide}
  The two sides of the symmetric interference pattern
\end{definition}
\begin{proof}
  This is the declared model definition of Fringe Side.
\end{proof}
\begin{definition}[sign]
  \label{def:physics:phyx-mini-0248:phyxminiproblems-problemphyxmini0248-fringeside-sign}
  \lean{PhyXMiniProblems.ProblemPhyXMini0248.FringeSide.sign}
  \uses{def:physics:phyx-mini-0248:phyxminiproblems-problemphyxmini0248-fringeside}
  Sign of a transverse screen coordinate on either side of the center
\end{definition}
\begin{proof}
  This is the declared model definition of sign.
\end{proof}
\begin{definition}[Fringe Kind]
  \label{def:physics:phyx-mini-0248:phyxminiproblems-problemphyxmini0248-fringekind}
  \lean{PhyXMiniProblems.ProblemPhyXMini0248.FringeKind}
  Bright maxima and dark minima in the screen pattern
\end{definition}
\begin{proof}
  This is the declared model definition of Fringe Kind.
\end{proof}
\begin{definition}[order Shift]
  \label{def:physics:phyx-mini-0248:phyxminiproblems-problemphyxmini0248-fringekind-ordershift}
  \lean{PhyXMiniProblems.ProblemPhyXMini0248.FringeKind.orderShift}
  \uses{def:physics:phyx-mini-0248:phyxminiproblems-problemphyxmini0248-fringekind}
  Offset from an integral path-difference multiple: a bright order m uses m, while a dark order m uses m + 1/2. Thus the question's order-two dark fringes have path-difference magnitude (2 + 1/2) λ
\end{definition}
\begin{proof}
  This is the declared model definition of order Shift.
\end{proof}
\begin{definition}[Light Source Kind]
  \label{def:physics:phyx-mini-0248:phyxminiproblems-problemphyxmini0248-lightsourcekind}
  \lean{PhyXMiniProblems.ProblemPhyXMini0248.LightSourceKind}
  The source technology named in the question
\end{definition}
\begin{proof}
  This is the declared model definition of Light Source Kind.
\end{proof}
\begin{definition}[Incident Wavefront]
  \label{def:physics:phyx-mini-0248:phyxminiproblems-problemphyxmini0248-incidentwavefront}
  \lean{PhyXMiniProblems.ProblemPhyXMini0248.IncidentWavefront}
  Wavefront shape incident on the double slit
\end{definition}
\begin{proof}
  This is the declared model definition of Incident Wavefront.
\end{proof}
\begin{definition}[Aperture Model]
  \label{def:physics:phyx-mini-0248:phyxminiproblems-problemphyxmini0248-aperturemodel}
  \lean{PhyXMiniProblems.ProblemPhyXMini0248.ApertureModel}
  Each sufficiently narrow opening is modeled as a point-like wave source
\end{definition}
\begin{proof}
  This is the declared model definition of Aperture Model.
\end{proof}
\begin{definition}[Figure Stage]
  \label{def:physics:phyx-mini-0248:phyxminiproblems-problemphyxmini0248-figurestage}
  \lean{PhyXMiniProblems.ProblemPhyXMini0248.FigureStage}
  The four explanatory stages explicitly shown in the supplied diagram
\end{definition}
\begin{proof}
  This is the declared model definition of Figure Stage.
\end{proof}
\begin{definition}[Double Slit Interference Setup]
  \label{def:physics:phyx-mini-0248:phyxminiproblems-problemphyxmini0248-doubleslitinterferencesetup}
  \lean{PhyXMiniProblems.ProblemPhyXMini0248.DoubleSlitInterferenceSetup}
  \uses{def:physics:phyx-mini-0248:phyxminiproblems-problemphyxmini0248-lengthquantity, def:physics:phyx-mini-0248:phyxminiproblems-problemphyxmini0248-signedlengthquantity, def:physics:phyx-mini-0248:phyxminiproblems-problemphyxmini0248-slitlabel, def:physics:phyx-mini-0248:phyxminiproblems-problemphyxmini0248-fringeside, def:physics:phyx-mini-0248:phyxminiproblems-problemphyxmini0248-fringekind, def:physics:phyx-mini-0248:phyxminiproblems-problemphyxmini0248-lightsourcekind, def:physics:phyx-mini-0248:phyxminiproblems-problemphyxmini0248-incidentwavefront, def:physics:phyx-mini-0248:phyxminiproblems-problemphyxmini0248-aperturemodel, def:physics:phyx-mini-0248:phyxminiproblems-problemphyxmini0248-figurestage}
  Physical data and observables of the double-slit experiment. fringeCoordinate kind side m is the signed transverse screen position of a labeled fringe. No field assigns the requested order-two separation or an answer-choice value
\end{definition}
\begin{proof}
  This is the declared model definition of Double Slit Interference Setup.
\end{proof}
\begin{definition}[Matches Scenario And Figure]
  \label{def:physics:phyx-mini-0248:phyxminiproblems-problemphyxmini0248-matchesscenarioandfigure}
  \lean{PhyXMiniProblems.ProblemPhyXMini0248.MatchesScenarioAndFigure}
  \uses{def:physics:phyx-mini-0248:phyxminiproblems-problemphyxmini0248-signedlengthreadout, def:physics:phyx-mini-0248:phyxminiproblems-problemphyxmini0248-slitlabel, def:physics:phyx-mini-0248:phyxminiproblems-problemphyxmini0248-fringeside, def:physics:phyx-mini-0248:phyxminiproblems-problemphyxmini0248-figurestage, def:physics:phyx-mini-0248:phyxminiproblems-problemphyxmini0248-doubleslitinterferencesetup}
  Qualitative information supplied by the prose and primary figure. The same normally incident plane wave illuminates two point-like narrow coherent apertures in phase. The figure labels the central maximum as m = 0 and the subsequent bright maxima on both sides by m = 1, 2, 3, 4. This predicate contains no numerical dark-fringe separation
\end{definition}
\begin{proof}
  This is the declared model definition of Matches Scenario And Figure.
\end{proof}
\begin{definition}[Matches Problem Readouts]
  \label{def:physics:phyx-mini-0248:phyxminiproblems-problemphyxmini0248-matchesproblemreadouts}
  \lean{PhyXMiniProblems.ProblemPhyXMini0248.MatchesProblemReadouts}
  \uses{def:physics:phyx-mini-0248:phyxminiproblems-problemphyxmini0248-lengthinmeters, def:physics:phyx-mini-0248:phyxminiproblems-problemphyxmini0248-lengthinmillimeters, def:physics:phyx-mini-0248:phyxminiproblems-problemphyxmini0248-lengthinnanometers, def:physics:phyx-mini-0248:phyxminiproblems-problemphyxmini0248-doubleslitinterferencesetup}
  Numerical readouts stated in the question: wavelength 633 nm, slit spacing 0.40 mm, and slit-to-screen distance 2.0 m
\end{definition}
\begin{proof}
  This is the declared model definition of Matches Problem Readouts.
\end{proof}
\begin{definition}[Has Physical Parameters]
  \label{def:physics:phyx-mini-0248:phyxminiproblems-problemphyxmini0248-hasphysicalparameters}
  \lean{PhyXMiniProblems.ProblemPhyXMini0248.HasPhysicalParameters}
  \uses{def:physics:phyx-mini-0248:phyxminiproblems-problemphyxmini0248-lengthinmeters, def:physics:phyx-mini-0248:phyxminiproblems-problemphyxmini0248-doubleslitinterferencesetup}
  Positive, nondegenerate optical and geometric parameters
\end{definition}
\begin{proof}
  This is the declared model definition of Has Physical Parameters.
\end{proof}
\begin{definition}[path Length From Slit Readout]
  \label{def:physics:phyx-mini-0248:phyxminiproblems-problemphyxmini0248-pathlengthfromslitreadout}
  \lean{PhyXMiniProblems.ProblemPhyXMini0248.pathLengthFromSlitReadout}
  \uses{def:physics:phyx-mini-0248:phyxminiproblems-problemphyxmini0248-signedlengthquantity, def:physics:phyx-mini-0248:phyxminiproblems-problemphyxmini0248-lengthreadout, def:physics:phyx-mini-0248:phyxminiproblems-problemphyxmini0248-signedlengthreadout, def:physics:phyx-mini-0248:phyxminiproblems-problemphyxmini0248-slitlabel, def:physics:phyx-mini-0248:phyxminiproblems-problemphyxmini0248-doubleslitinterferencesetup}
  Euclidean path length from one slit to a screen coordinate, read in one common length unit. The slit plane is at axial coordinate zero, the screen is at screenDistance, and the slit coordinates are ± slitSeparation / 2
\end{definition}
\begin{proof}
  This is the declared model definition of path Length From Slit Readout.
\end{proof}
\begin{definition}[Is Finite Fringe Order]
  \label{def:physics:phyx-mini-0248:phyxminiproblems-problemphyxmini0248-isfinitefringeorder}
  \lean{PhyXMiniProblems.ProblemPhyXMini0248.IsFiniteFringeOrder}
  \uses{def:physics:phyx-mini-0248:phyxminiproblems-problemphyxmini0248-lengthinmeters, def:physics:phyx-mini-0248:phyxminiproblems-problemphyxmini0248-fringekind, def:physics:phyx-mini-0248:phyxminiproblems-problemphyxmini0248-doubleslitinterferencesetup}
  Only path differences smaller than the slit spacing correspond to a fringe at a finite screen coordinate in the exact two-point-source geometry. The condition is stated in meters, but is independent of the selected unit
\end{definition}
\begin{proof}
  This is the declared model definition of Is Finite Fringe Order.
\end{proof}
\begin{definition}[Satisfies Exact Two Point Source Laws]
  \label{def:physics:phyx-mini-0248:phyxminiproblems-problemphyxmini0248-satisfiesexacttwopointsourcelaws}
  \lean{PhyXMiniProblems.ProblemPhyXMini0248.SatisfiesExactTwoPointSourceLaws}
  \uses{def:physics:phyx-mini-0248:phyxminiproblems-problemphyxmini0248-lengthreadout, def:physics:phyx-mini-0248:phyxminiproblems-problemphyxmini0248-signedlengthreadout, def:physics:phyx-mini-0248:phyxminiproblems-problemphyxmini0248-fringeside, def:physics:phyx-mini-0248:phyxminiproblems-problemphyxmini0248-fringekind, def:physics:phyx-mini-0248:phyxminiproblems-problemphyxmini0248-doubleslitinterferencesetup, def:physics:phyx-mini-0248:phyxminiproblems-problemphyxmini0248-pathlengthfromslitreadout, def:physics:phyx-mini-0248:phyxminiproblems-problemphyxmini0248-isfinitefringeorder}
  Exact scalar-wave interference laws for two in-phase point-like sources and a finite-distance planar screen. For every physically finite fringe, the difference between the Euclidean path lengths from the lower and upper slits is the appropriate integer or half-integer multiple of the wavelength. The second field records which side of the optical axis the named coordinate lies on. Neither field mentions the distance between the two order-two dark fringes or any answer choice
\end{definition}
\begin{proof}
  This is the declared model definition of Satisfies Exact Two Point Source Laws.
\end{proof}
\begin{definition}[dark Fringe Separation In Millimeters]
  \label{def:physics:phyx-mini-0248:phyxminiproblems-problemphyxmini0248-darkfringeseparationinmillimeters}
  \lean{PhyXMiniProblems.ProblemPhyXMini0248.darkFringeSeparationInMillimeters}
  \uses{def:physics:phyx-mini-0248:phyxminiproblems-problemphyxmini0248-signedlengthreadout, def:physics:phyx-mini-0248:phyxminiproblems-problemphyxmini0248-doubleslitinterferencesetup}
  Distance between the two symmetric dark fringes of a given order, read in millimeters. This is a generic observable derived from their signed screen coordinates, not a definition of the numerical answer
\end{definition}
\begin{proof}
  This is the declared model definition of dark Fringe Separation In Millimeters.
\end{proof}
\begin{definition}[finite Screen Symmetric Separation]
  \label{def:physics:phyx-mini-0248:phyxminiproblems-problemphyxmini0248-finitescreensymmetricseparation}
  \lean{PhyXMiniProblems.ProblemPhyXMini0248.finiteScreenSymmetricSeparation}
  Exact separation of the symmetric solutions for axial distance L, slit spacing d, and positive path-difference magnitude δ in a common unit. This generic geometric expression is not specialized to the problem data
\end{definition}
\begin{proof}
  This is the declared model definition of finite Screen Symmetric Separation.
\end{proof}
\begin{definition}[Answer Choice]
  \label{def:physics:phyx-mini-0248:phyxminiproblems-problemphyxmini0248-answerchoice}
  \lean{PhyXMiniProblems.ProblemPhyXMini0248.AnswerChoice}
  The four answer labels printed with the problem
\end{definition}
\begin{proof}
  This is the declared model definition of Answer Choice.
\end{proof}
\begin{definition}[displayed Distance In Millimeters]
  \label{def:physics:phyx-mini-0248:phyxminiproblems-problemphyxmini0248-displayeddistanceinmillimeters}
  \lean{PhyXMiniProblems.ProblemPhyXMini0248.displayedDistanceInMillimeters}
  \uses{def:physics:phyx-mini-0248:phyxminiproblems-problemphyxmini0248-answerchoice}
  The distance in millimeters printed next to each answer label
\end{definition}
\begin{proof}
  This is the declared model definition of displayed Distance In Millimeters.
\end{proof}
\begin{definition}[recorded Dataset Answer]
  \label{def:physics:phyx-mini-0248:phyxminiproblems-problemphyxmini0248-recordeddatasetanswer}
  \lean{PhyXMiniProblems.ProblemPhyXMini0248.recordedDatasetAnswer}
  \uses{def:physics:phyx-mini-0248:phyxminiproblems-problemphyxmini0248-answerchoice}
  The answer label recorded by the source dataset
\end{definition}
\begin{proof}
  This is the declared model definition of recorded Dataset Answer.
\end{proof}
\begin{definition}[rounded To Nearest Tenth Millimeter]
  \label{def:physics:phyx-mini-0248:phyxminiproblems-problemphyxmini0248-roundedtonearesttenthmillimeter}
  \lean{PhyXMiniProblems.ProblemPhyXMini0248.roundedToNearestTenthMillimeter}
  Round a millimeter readout to the nearest tenth of a millimeter
\end{definition}
\begin{proof}
  This is the declared model definition of rounded To Nearest Tenth Millimeter.
\end{proof}
\begin{definition}[Matches Displayed Distance]
  \label{def:physics:phyx-mini-0248:phyxminiproblems-problemphyxmini0248-matchesdisplayeddistance}
  \lean{PhyXMiniProblems.ProblemPhyXMini0248.MatchesDisplayedDistance}
  \uses{def:physics:phyx-mini-0248:phyxminiproblems-problemphyxmini0248-doubleslitinterferencesetup, def:physics:phyx-mini-0248:phyxminiproblems-problemphyxmini0248-darkfringeseparationinmillimeters, def:physics:phyx-mini-0248:phyxminiproblems-problemphyxmini0248-answerchoice, def:physics:phyx-mini-0248:phyxminiproblems-problemphyxmini0248-displayeddistanceinmillimeters, def:physics:phyx-mini-0248:phyxminiproblems-problemphyxmini0248-roundedtonearesttenthmillimeter}
  A displayed choice matches the physical result when the exact finite-screen separation, rounded to the precision used by all four choices, equals its printed value
\end{definition}
\begin{proof}
  This is the declared model definition of Matches Displayed Distance.
\end{proof}
\begin{definition}[Is Unique Matching Displayed Distance]
  \label{def:physics:phyx-mini-0248:phyxminiproblems-problemphyxmini0248-isuniquematchingdisplayeddistance}
  \lean{PhyXMiniProblems.ProblemPhyXMini0248.IsUniqueMatchingDisplayedDistance}
  \uses{def:physics:phyx-mini-0248:phyxminiproblems-problemphyxmini0248-doubleslitinterferencesetup, def:physics:phyx-mini-0248:phyxminiproblems-problemphyxmini0248-answerchoice, def:physics:phyx-mini-0248:phyxminiproblems-problemphyxmini0248-matchesdisplayeddistance}
  A choice is the unique displayed value matching the physical result
\end{definition}
\begin{proof}
  This is the declared model definition of Is Unique Matching Displayed Distance.
\end{proof}
\begin{lemma}[order Two Dark Fringe Separation exact Finite Screen]
  \label{lem:physics:phyx-mini-0248:phyxminiproblems-problemphyxmini0248-ordertwodarkfringeseparation-exactfinitescreen}
  \lean{PhyXMiniProblems.ProblemPhyXMini0248.orderTwoDarkFringeSeparation_exactFiniteScreen}
  \uses{def:physics:phyx-mini-0248:phyxminiproblems-problemphyxmini0248-doubleslitinterferencesetup, def:physics:phyx-mini-0248:phyxminiproblems-problemphyxmini0248-matchesproblemreadouts, def:physics:phyx-mini-0248:phyxminiproblems-problemphyxmini0248-hasphysicalparameters, def:physics:phyx-mini-0248:phyxminiproblems-problemphyxmini0248-satisfiesexacttwopointsourcelaws, def:physics:phyx-mini-0248:phyxminiproblems-problemphyxmini0248-darkfringeseparationinmillimeters, def:physics:phyx-mini-0248:phyxminiproblems-problemphyxmini0248-finitescreensymmetricseparation}
  The exact finite-screen geometry gives the radical expression below. Here 633 / 400000 mm is the order-two dark path difference (2 + 1/2) · 633 nm. This is approximately 15.8251239 mm, rather than identically the paraxial value 633 / 40 mm = 15.825 mm
\end{lemma}
\begin{proof}
  The conclusion follows from the governing relations and figure readouts named in the statement of order Two Dark Fringe Separation exact Finite Screen.
\end{proof}
% --- Archon declaration topology end ---
% --- Archon physics formalization source end ---
